\chapter{Linear Equations in One Variable}

\chapterintro{Solving an equation means undoing, in reverse order, everything that was done to the variable, while keeping both sides equal. Checking the answer by substitution is the first unit test you will ever write by hand.}

\section{Core Ideas}

\paragraph{The balance rule} An equation stays true if you do the same thing to both sides: add, subtract, multiply or divide by the same nonzero number. Inverse operations undo each other: subtraction undoes addition, division undoes multiplication.

\paragraph{A solving procedure} Simplify each side (distribute, combine like terms). Move variable terms to one side and constants to the other. Divide by the coefficient of the variable. Check by substituting into the original equation.

\paragraph{Three outcomes} Most equations have one solution. If the variables cancel and you are left with a false statement such as $0 = 7$, there is no solution. If you are left with a true statement such as $0 = 0$, every number is a solution (an identity). Your solver in Exercise~C3.1 must return all three.

\paragraph{Literal equations} Solving $P = 2l + 2w$ for $w$ uses the same moves and gives $w = (P - 2l)/2$. This is how you turn one formula into its inverse, which you will do constantly when converting units or reversing a transformation.

\paragraph{Proportions} $a/b = c/d$ is equivalent to $ad = bc$ (cross-multiplication), provided $b$ and $d$ are nonzero.

\paragraph{Word problems} Name the unknown with its unit, write the relationship as an equation, solve, then check the answer against the story, not just the equation.

\section{Worked Examples}

\begin{example}
Solve $3x - 7 = 11$.
\solution Add 7: $3x = 18$. Divide by 3: $x = 6$. Check: $3(6) - 7 = 11$.
\end{example}

\begin{example}
Solve $2(x + 3) = 5x - 9$.
\solution Distribute: $2x + 6 = 5x - 9$. Subtract $2x$: $6 = 3x - 9$. Add 9: $15 = 3x$. So $x = 5$. Check: $2(8) = 16$ and $25 - 9 = 16$.
\end{example}

\begin{example}
Solve $A = \tfrac{1}{2}bh$ for $h$.
\solution Multiply by 2: $2A = bh$. Divide by $b$: $h = 2A/b$.
\end{example}

\section{Exercises}

\begin{multicols}{2}
\begin{exercises}
\item Solve $5x + 4 = 29$.
\item Solve $\dfrac{x}{3} - 2 = 4$.
\item Solve $4(x - 2) = 2x + 10$.
\item Solve $7 - 2(x + 1) = 3x - 5$.
\item Solve $3(2x - 1) = 6x + 4$. State what kind of outcome this is.
\item Solve $2(x + 4) - x = x + 8$. State what kind of outcome this is.
\item Solve $F = \tfrac{9}{5}C + 32$ for $C$.
\item Solve $\dfrac{3}{x} = \dfrac{12}{20}$.
\item Solve $\dfrac{x + 1}{4} = \dfrac{2x - 3}{2}$. Give the exact answer as a fraction.
\item A consultant charges a \$150 setup fee plus \$85 per hour. A client was billed \$830. How many hours were worked?
\item Two trains start 300 miles apart and travel toward each other, one at 60 mph and the other at 40 mph. After how many hours do they meet?
\end{exercises}
\end{multicols}

\section{Programming Exercises}

\begin{codeexercises}
\item Write \lstinline!solve_linear(a, b, c, d)! that solves $ax + b = cx + d$ and returns one of: \lstinline!("one", x)!, \lstinline!("none",)! or \lstinline!("all",)!. Test it against Exercises~3.1 through 3.6.
\item Write \lstinline!check_linear(a, b, c, d, x)! that substitutes $x$ into both sides and compares them with \lstinline!math.isclose!. Use it to verify every \lstinline!("one", x)! result from C3.1, so your solver and your checker test each other.
\item Write \lstinline!c_to_f(c)! and \lstinline!f_to_c(f)! from Exercise~3.7. Assert that \lstinline!f_to_c(c_to_f(c))! returns \lstinline!c! (within tolerance) for every integer $c$ from $-50$ to $150$.
\item Rewrite C3.1 to use \lstinline!fractions.Fraction! for all arithmetic so that Exercise~3.9 returns exactly $7/3$ rather than $2.3333$. Decide what the function should do if the inputs are floats.
\end{codeexercises}
